\section{Comparison with Topmetal-S}
\label{app:topmetal}

Topmetal-S and NAUSICA use exposed metal electrodes connected to CSAs to
measure charge without gas amplification. A focusing electrode guides the
ions onto the collecting pads. Topmetal-S was developed to measure both
energy and track topology in a high-pressure ion TPC~\cite{mei2020topmetalcmosdirectcharge}.
In ITACA, electroluminescence provides the energy measurement, allowing
NAUSICA to choose its pitch and readout for ion-track imaging.

The Topmetal-S collecting electrode is a hexagon of \SI{1}{mm} diameter,
implemented in a 130 nm CMOS process~\cite{Liang2024TopmetalS}. The
N$\nu$DEx readout model places these electrodes at \SI{10}{mm}
pitch~\cite{Liang2026NvDExSimulation}. NAUSICA uses \SI{0.5}{mm} square
pads at \SI{2}{mm} pitch, close to the diffusion scale in ITACA. This
finer pitch samples the track with more pixels. The smaller pad reduces
the capacitive load, while the focusing field provides collection over
the pixel cell. The charge measured in each pixel must preserve the
track and its blob structure; the precise total-energy measurement is
provided by the electron signal.

\subsection{Noise measured with charge injection}

Topmetal-S reports an ENC of approximately $115\,e^-$ after a digital
trapezoidal filter~\cite{Yang2024TopmetalCSA}. Repeated injections of a
fixed charge produce a distribution of filtered pulse amplitudes. For
an injected charge of $N_{\rm inj}$ elementary charges, the ENC is
\begin{equation}
 \mathrm{ENC}=N_{\rm inj}\frac{\sigma_A}{\langle A\rangle},
\end{equation}
where $\langle A\rangle$ and $\sigma_A$ are the mean and standard deviation
of that distribution. The quoted noise therefore includes the effect
of the filter.

We perform a corresponding injection calculation for the circuit in
Table~\ref{tab:spice_devices}, with an external capacitance of
\SI{\SpicePadPf}{pF}. A pulse carrying
$\SpiceInjectionCalibration\,e$ has a \SI{1}{\micro\second} plateau
and \SI{10}{ns} rise and fall times. We apply a trapezoidal filter with
\SI{300}{\micro\second} rise time and \SI{300}{\micro\second} flat top.
The filtered voltage per elementary charge calibrates the signal. Applying
the same filter to the SPICE noise spectrum, including thermal and flicker
noise over \SI{\SpiceLowHz}{Hz}--\SI{\SpiceHighMHz}{MHz}, gives an ENC of
$\SpiceInjectionEnc\,e^-$ RMS. The calculation uses the supply and
temperature in Table~\ref{tab:spice_parameters} and noise evaluated at
the DC operating point.

\begin{table}[!htbp]
 \centering
 \begin{tabular}{lll}
 \toprule
 & Topmetal-S & NAUSICA \\
 \midrule
 CMOS process & 130 nm & 180 nm \\
 Collecting electrode & \SI{1}{mm} diameter, hexagonal & \SI{0.5}{mm} side, square \\
 Pitch & \SI{10}{mm} (N$\nu$DEx model) & \SI{2}{mm} \\
 Injection ENC & $\sim115\,e^-$, measured & $\SpiceInjectionEnc\,e^-$, simulated \\
 \bottomrule
 \end{tabular}
 \caption{Geometry and charge-injection noise comparison. Both ENC values
 include trapezoidal filtering. The NAUSICA filter settings are given above;
 the Topmetal-S report~\cite{Yang2024TopmetalCSA} does not specify its
 trapezoid rise and flat-top durations.}
 \label{tab:topmetal_comparison}
\end{table}
